\uuid{6act}
\titre{Probabilité qu'une variable normale dépasse une autre}
\auteur{Maxime Nguyen}
\organisation{AMSCC}
\datecreate{2026-11-27}
\niveau{L2}
\module{Probabilité}
\chapitre{Probabilité continue}
\sousChapitre{Loi normale}
\difficulte{2}
\theme{loi normale, différence de variables indépendantes, standardisation}

\contenu{
\begin{question}
\qtype{onlyone}
On note $\mathcal{N}(\mu,\sigma)$ la loi normale d'espérance $\mu$ et d'écart-type $\sigma$, et $\Phi$ la fonction de répartition de $\mathcal{N}(0,1)$. Soit $X \sim \mathcal{N}(10, 3)$ et $Y \sim \mathcal{N}(6, 4)$, indépendantes. $\mathbb{P}(X > Y)$ vaut :
\begin{answers}
\bad{$1 - \Phi(0{,}8)$}
\bad{$\Phi(4)$}
\bad{$\Phi\left(\frac47\right)$}
\good{$\Phi(0{,}8)$}
\end{answers}
\begin{explanations}
$D = X - Y$ est normale, d'espérance $10 - 6 = 4$ et de
variance $9 + 16 = 25$ (les variances s'ajoutent, même pour une
différence), donc $D \sim \mathcal{N}(4, 5)$. Puis
$\mathbb{P}(D > 0) = \mathbb{P}\left(Z > -\frac45\right) = 1 - \Phi(-0{,}8) = \Phi(0{,}8)$.

A : erreur de signe au moment de la symétrie ; contrôle : $\mathbb{E}(D) > 0$, donc $\mathbb{P}(D > 0) > \frac12$.

B : on a soustrait les écarts-types ($4 - 3 = 1$).

C : on a ajouté les écarts-types ($3 + 4 = 7$).

Contre-sens visé : croire que la dispersion d'une différence est la
différence des dispersions.
\end{explanations}
\end{question}
}
